\begin{table}[t]
\centering
\caption{Marine protected areas in the study, 2019--2025 (hours of apparent fishing effort).}\label{tab:mpas}
\small
\begin{tabular}{lrrrrrrrr}
\toprule
MPA & region & type & cells & interior & y\_nt & y\_inside & y\_press & share (\%) \\
\midrule
Nazca-Desventuradas & oceanic & no-take & 6,700 & 665 & 138 & 138 & 1,962,168 & 56.3 \\
Rapa Nui & oceanic & multi-use & 11,535 & 1,307 & 0 & 162 & 1,011,870 & 29.0 \\
Montes submarinos Crusoe y Selkirk & oceanic & no-take & 283 & 3 & 0 & 0 & 110,237 & 3.2 \\
Kawésqar (RN) & coastal\_south & multi-use & 1,928 & 314 & 0 & 1,447 & 99,887 & 2.9 \\
Tortel & coastal\_south & multi-use & 865 & 78 & 0 & 5,332 & 63,451 & 1.8 \\
Lafken Mapu Lahual & coastal\_south & multi-use & 444 & 0 & 0 & 0 & 45,896 & 1.3 \\
Islas Diego Ramírez y Paso Drake & coastal\_south & no-take & 6,169 & 2,187 & 1,656 & 1,656 & 40,958 & 1.2 \\
Tictoc-Golfo Corcovado & coastal\_south & no-take & 348 & 12 & 27 & 27 & 37,696 & 1.1 \\
Bahía Lomas & coastal\_south & multi-use & 463 & 7 & 0 & 0 & 34,542 & 1.0 \\
Mar de Pisagua & coastal\_north & multi-use & 363 & 6 & 0 & 794 & 27,669 & 0.8 \\
Mar de Juan Fernández (PN) & oceanic & no-take & 2,922 & 643 & 23 & 23 & 25,915 & 0.7 \\
Estero de Quitralco & coastal\_south & multi-use & 147 & 0 & 0 & 0 & 12,488 & 0.4 \\
Mar de Juan Fernández & oceanic & multi-use & 746 & 58 & 0 & 7 & 5,928 & 0.2 \\
Punta Morro - Desembocadura Río Copiapó - Isla grande de Atacama & coastal\_north & multi-use & 297 & 1 & 0 & 13 & 4,581 & 0.1 \\
Pitipalena - Añihue & coastal\_south & multi-use & 62 & 1 & 0 & 0 & 1,775 & 0.1 \\
Seno Almirantazgo & coastal\_south & multi-use & 216 & 7 & 0 & 0 & 1,481 & 0.0 \\
Archipiélago de Humboldt & coastal\_north & multi-use & 612 & 51 & 0 & 112 & 1,236 & 0.0 \\
Francisco Coloane (AM) & coastal\_south & multi-use & 16 & 9 & 0 & 16 & 3 & 0.0 \\
Isla Chañaral & coastal\_north & multi-use & 1 & 1 & 0 & 1 & 0 & 0.0 \\
Motu Motiro Hiva & oceanic & no-take & 344 & 344 & 1 & 1 & 0 & 0.0 \\
\bottomrule
\end{tabular}
\end{table}